\documentclass[11pt]{article}
\usepackage[a4paper,margin=28mm]{geometry}
\usepackage[T1]{fontenc}
\usepackage{lmodern}
\usepackage{amsmath,amssymb}
\emergencystretch=2em
\title{The identity operator disproves the half upper bound\\Conjecture 00000007167}
\author{ziangni-sys}
\date{4 October 2026}
\begin{document}
\maketitle
\section*{The original assertion}
The English source concerns equivalence of numerical radius and operator
norm and asserts that the upper bound of the ratio of numerical radius
to operator norm is one half. The Chinese source states the same direction
of the ratio and the same upper bound.
The additional claim about tightness at a two-dimensional rotation is
part of the conjunction; disproving the universal upper bound already
disproves the conjecture.

\section*{An actual bounded linear operator}
Let $H=\mathbb C$ with its ordinary Hilbert norm and inner product, and
let $I:H\to H$ be the identity operator. Define
\[
\|A\|=\sup_{\|v\|=1}\|Av\|,
\qquad
w(A)=\sup_{\|v\|=1}|\langle v,Av\rangle|.
\]
The identity is complex-linear and bounded. Every unit vector has
$\|Iv\|=\|v\|=1$, and unit vectors exist, for example $v=1$.
Thus $\|I\|=1$.
For every unit vector,
\[
|\langle v,Iv\rangle|=\langle v,v\rangle=\|v\|^2=1.
\]
The set of real values whose supremum defines $w(I)$ is exactly
the singleton $\{1\}$, so $w(I)=1$. Consequently
\[
\frac{w(I)}{\|I\|}=1>\frac12.
\]
This violates the proposed bound. The same calculation works for the
identity on any nonzero Hilbert space, including a two-dimensional one.
The numerical radius here is not a spectral-radius surrogate: its full
unit-vector supremum is computed directly.

\section*{Formal verification}
The Lean project uses the standard complex inner product and norm, and
the actual bounded operator
\texttt{ContinuousLinearMap.id} on $\mathbb C$.
The operator-norm calculation uses Mathlib's theorem for the norm of
the identity.
The numerical-value set is defined by existence of a complex unit vector
with the corresponding norm of its inner product.
The proof establishes equality of this entire set with $\{1\}$;
the radius is its real supremum, rather than an assumed constant.
The final theorem \texttt{conjecture\_7167\_false} refutes the claimed
upper bound for all positive-norm operators on this Hilbert space.
Any universal upper bound as in the original must apply to this case.
There are no auxiliary computations or unproved analytic estimates.
\end{document}
